% Contenido exclusivo del propietario de entropías tipadas.
\subsection{Quasi-free correlation and spectral occupation entropies}

For a gauge-invariant, number-conserving quasifree fermionic state
defined on a finite-dimensional one-particle space, let
\(C_{ij}=\langle a_j^\dagger a_i\rangle\) be the one-particle correlation
operator. The gauge-invariant hypothesis excludes anomalous pairing
correlations, so that \(C\) determines the quasifree state. If
\(0<C<I\), its one-particle modular generator and the reciprocal reconstruction
of the correlator are
\begin{equation}
 \boxed{
 K_1=\log\bigl((I-C)C^{-1}\bigr),
 \qquad
 C=(I+e^{K_1})^{-1}.}
 \label{eq:puente-cuasilibre-modular-rev3}
\end{equation}

\begin{corollary}[Logistic form of the modular spectrum]
\label{cor:forma-logistica-modular-rev3}
If \(K_1u_i=\kappa_i u_i\), the occupation of mode \(u_i\) is
\[
 \langle n_i\rangle=\frac1{1+e^{\kappa_i}}.
\]
For \(\kappa_i=\beta(\varepsilon_i-\mu)\), the Fermi--Dirac law
is recovered.
\end{corollary}

\begin{proof}
The spectral diagonalization of \(K_1\) diagonalizes every function thereof. Applying
\(x\mapsto(1+e^x)^{-1}\) to each eigenvalue yields the formula.
\end{proof}

If \(C\) has eigenvalues \(p_i\in[0,1]\), the occupation entropy of the
quasifree fermionic state is, with the convention \(0\log0=0\),
\begin{equation}
 S_F=-\sum_i\bigl[p_i\log p_i+(1-p_i)\log(1-p_i)\bigr].
 \label{eq:entropia-cuasilibre-fermi-rev3}
\end{equation}
For a diagonal Gaussian bosonic state with finitely many modes and mean
occupancies \(\overline n_i\),
\begin{equation}
 S_B=\sum_i\bigl[(1+\overline n_i)\log(1+\overline n_i)
                  -\overline n_i\log\overline n_i\bigr].
 \label{eq:entropia-gaussiana-bose-rev3}
\end{equation}
The same formula extends to a countable family when the series on the
right-hand side converges.
These entropies quantify occupation configurations. Entanglement entropy
has another domain: the reduced state induced by a tensor
decomposition. Identifying the two requires specifying the reduction and the
correlation operator that implements it.

Finally, for a reduced state \(\rho_A\), the \emph{von Neumann
entropy}
\[
 S_{\rm vN}(\rho_A)=-\operatorname{Tr}(\rho_A\log\rho_A)
\]
is a spectral property of the density operator. When \(\rho_A\)
comes from a global purification, it quantifies entanglement across
the cut. In the HMT boundary construction, the area operator, the
modular Hamiltonian, and the corresponding reduction remain in the
\cref{hap:sec:area-entrelazamiento-barbero-rev11}.

The maps between these domains are published through explicit identities.
The decomposition of a measure along the fibres of \(q\) satisfies the chain
rule
\begin{equation}
 \mathsf H_\mu(\mathbf X)
 =\mathsf H_{q_*\mu}(q(\mathbf X))
  +\mathsf H_{\rm fib}(q;\mu).
 \label{eq:cadena-entropia-fibras-rev2}
\end{equation}
\Needspace{5\baselineskip}
A finite distribution \(p=(p_i)\) is realized diagonally as
\[
 \rho_p=\sum_i p_i|i\rangle\langle i|,
 \qquad S_{\rm vN}(\rho_p)=\mathsf H(p).
\]
If \(P_{N,E}\) projects onto the microcanonical subspace of dimension
\(\Omega(N,E)\), then
\[
 \rho_{\rm mc}=\frac{P_{N,E}}{\Omega(N,E)},
 \qquad S_{\rm vN}(\rho_{\rm mc})=\log\Omega(N,E)=S_{\rm mc}.
\]
For a pure bipartite state,
\[
 \rho_A=\operatorname{Tr}_B
 \bigl(|\Psi_{AB}\rangle\langle\Psi_{AB}|\bigr)
\]
defines the reduction that measures entanglement. Finally, equality between
\(h_{\rm top}=\log\rho(A)\) and a metric entropy is attained, for this
primitive system, by means of the Parry measure, which realizes maximal metric
entropy; it is not attributed to an arbitrary measure on histories.

